\section{Conclusion}

%Existing scholarship on the role of landlords in housing inequality and segregation has been concerned with discrimination as a reason for segregation and inequality \citep{charlesDynamicsRacialResidential2003}. In particular, there have been two theories of discrimination that have been applied to housing as well \citep{auspurgClosedDoorsEverywhere2019}: taste-based and statistical discrimination. Taste-based discrimination refers to to a landlord/realtor that dismisses applications from households with sufficient funds based on their group membership and is willing to accept costs from dismissing potential customers. Economists have often seen this type of discrimination as inefficient and that market competition should make this type of discrimination less prevalent \citep{arrowWhatHasEconomics1998}. Another strand of theory, and a partial response to this observation, has been statistical discrimination where the landlords uses group membership as proxy. This theory has proven the possibility a third reason of apparent discrimination in housing which in existing tests cannot be distinguished from taste-based discrimination on behalf of landlords. But when households value their neighborhood, renting to minority households decreases local rents and makes investments less attractive. With lower investments into housing quality, minority neighborhoods will be of lower quality even if there is sufficient income. And the other way around, minority households need to pay a premium to achieve a similar residential outcomes to offset their effect on the neighborhood. Similar arguments have been made in labor markets where employers might fear that their existing workforce opposes hiring minority members \citep[95]{arrowWhatHasEconomics1998} and aligns with the observed concern among existing residents that housing values will decrease when minority members move to the neighborhood \citep{harrisPropertyValuesDrop1999, krysanDoesRaceMatter2009}. 

The proposed theoretical model extends existing models of residential segregation and housing inequalities. By incorporating the characteristics of housing units and the agency of landlords into neighborhood composition, the model extends beyond existing explanations that primarily emphasize demand-side factors \citep{desmondPoorPayMore2019, rosenthalChangePersistenceEconomic2015}. This addition is crucial, as empirical research suggests that households often prioritize features of their homes—such as quality and affordability—over the surrounding neighborhood's amenities or demographics. Ignoring the built environment's role in explaining residential segregation and neighborhood change misrepresents the motives and behaviors of households. The model shows that both households' neighborhood preferences and spatially contingent investments can create sizable levels of segregation (even without the presence of the other). When both mechanisms are at work, they can reinforce each other and create stable segregated neighborhoods by income, social status, as well as rent and housing quality. Moreover, the insights gained from this model reveal that the mechanisms underlying residential segregation and neighborhood change are intrinsically linked to patterns of residential mobility and housing inequality. By considering housing quality in segregation models, these models become more realistic (they better reflect household preferences) and more general (they explain a wider range of phenomena). This model, therefore, represents a crucial step toward a more comprehensive synthesis of urban theory.

%This model deliberately concentrated on socio-economic segregation with little reference to racial/ethnic segregation. But it could be extended similarly to the models by \citet{masseyAmericanApartheidSegregation1993, bruchHowPopulationStructure2014}, where ethnic inequality and economic segregation interact, explaining patterns in ethnic segregation. Although there is a large body of work that documents how race and ethnicity shape residential segregation and neighborhood change in the US (e.g., \cite{hwangUnequalDisplacementGentrification2020}), these results are likely not generalizable as the history of racial relations and institutionalization of racial boundaries in the US is very specific. The extent to which racial and ethnic composition plays a role in residential segregation in other contexts beyond economic disadvantage is underresearched. Residential choices are also connected to choices in other domains. For example, research has documented that schools are an important consideration of family decisions about neighborhoods in the US \citep{CITE}. \citet{fieldCoupledDynamicsNeighborhood2026} recently proposed a joint model of demographic change in schools and neighborhoods.


%\section{Implications for Policy}

%Ghettos and very disadvantaged nbs are a result of the general workings of thehousing market, therefore place-based policies will not work. to prevent extremes, one needs to combat segregation as a whole. but how? if characteristics of the dwelling are most important, one could try to foster mixed housing: mandate investors to build different sized and priced housing units in the same building in rich areas, programs to reduce investment uncertainty in poor areas, mixed zoning etc.

%In this theory, I take the position that neighborhood change is not a phenomenon on the neighborhood level but a manifestation of a shift in spatial equilibrium caused by external shocks like income or population increases (like e.g. \cite{galsterNatureNeighbourhood2001, leeNaturalAmenitiesNeighbourhood2018}). Neighborhood change is therefore caused by changes to the city as a whole but the workings of the housing market concentrate the resulting changes in space. I therefore disagree with \citet[358]{hwangUnequalDisplacementGentrification2020} who observe that "gentrification reduces the overall pool of neighborhoods accessible to financially disadvantaged residents within the city". Not gentrification "causes" less affordable housing, shocks in population and income cause it. Gentrification is the process by which it does. If we would try to protect individual neighborhoods, the increase in demand would simply create displacement elsewhere. Nonetheless, as moving includes transaction costs and displacement primarily affects the poor, there is indeed a social injustice connected to gentrification. But place-based policies will not work as demand will seek its supply somewhere else (\textit{substitution effect}). Gentrification should less be a concern in its own right but city-wide housing supply should, and particularly housing elasticity. 